\section{Planification du projet}

\subsection{Méthode de priorisation et d'estimation}

Deux informations distinctes sont associées à chaque user story :

\begin{itemize}[leftmargin=7mm]
  \item la \textbf{priorité} exprime la valeur métier et l'urgence. Les niveaux Haute,
  Moyenne et Faible correspondent respectivement à une adaptation des catégories
  \emph{Must}, \emph{Should} et \emph{Could} de MoSCoW ;
  \item les \textbf{story points} estiment l'effort relatif, l'incertitude et le risque
  selon la suite de Fibonacci : 1, 2, 3, 5, 8 et 13.
\end{itemize}

Les points ne représentent donc ni des jours ni une priorité. Une fonctionnalité
hautement prioritaire peut recevoir peu de points si elle est simple, tandis qu'une
évolution facultative comme le firmware peut recevoir 13 points en raison de son
risque technique.

\subsection{Product Backlog révisé}

Le backlog initial de 30 user stories a été enrichi à mesure que les relations
multi-organisations, la disponibilité calculée et les parcours réels ont été
clarifiés. La gestion des véhicules, retirée du produit après validation UX, est
remplacée par un onboarding adapté au rôle. Deux user stories distinguent désormais
la gestion commerciale globale de la consultation de l'abonnement par
l'administrateur. Le backlog révisé contient 42 user stories pour 301 points.

\begingroup
\footnotesize
\setlength{\tabcolsep}{3.5pt}
\begin{longtable}{|C{1.25cm}|p{10.25cm}|C{1.65cm}|C{1.0cm}|}
  \caption{Product Backlog de \ProjectName{}}
  \label{tab:product-backlog}\\
  \hline
  \textbf{ID} & \textbf{User Story} & \textbf{Priorité} & \textbf{Pts}\\
  \hline
  \endfirsthead
  \hline
  \textbf{ID} & \textbf{User Story} & \textbf{Priorité} & \textbf{Pts}\\
  \hline
  \endhead
  US-01 & En tant que visiteur, je veux créer un compte client ou me connecter de
  manière sécurisée par email ou Google. & Haute & 8\\\hline
  US-02 & En tant qu'utilisateur authentifié, je veux consulter et mettre à jour mon
  profil ainsi que les paramètres de sécurité de mon compte. & Moyenne & 3\\\hline
  US-03 & En tant que super administrateur, je veux gérer les organisations, leurs
  administrateurs et les accès globaux de la plateforme. & Haute & 8\\\hline
  US-04 & En tant qu'administrateur, je veux créer, modifier et désactiver les comptes
  des opérateurs et techniciens de mon organisation. & Haute & 5\\\hline
  US-05 & En tant que super administrateur, je veux consulter l'audit global, l'état
  des intégrations et les statistiques d'utilisation de la plateforme. & Moyenne & 5\\\hline
  US-06 & En tant qu'administrateur, je veux consulter les clients ayant utilisé les
  bornes de mon organisation avec leurs sessions et statistiques. & Moyenne & 5\\\hline
  US-07 & En tant qu'opérateur, je veux ajouter, modifier et désactiver les bornes de
  mon organisation avec leurs informations techniques et géographiques. & Haute & 8\\\hline
  US-08 & En tant qu'opérateur, je veux gérer les connecteurs associés aux bornes de
  mon organisation avec leur type, puissance et état. & Haute & 5\\\hline
  US-09 & En tant qu'opérateur, je veux consulter la liste des bornes avec leur
  disponibilité calculée, leur connexion et leur dernière activité. & Haute & 8\\\hline
  US-10 & En tant qu'opérateur, je veux visualiser les bornes sur une carte et ajouter
  une borne en sélectionnant son emplacement. & Haute & 8\\\hline
  US-11 & En tant que technicien, je veux consulter sans modification les bornes et
  leur emplacement sur la carte de mon organisation. & Haute & 5\\\hline
  US-12 & En tant que client, je veux rechercher et filtrer les bornes publiques selon
  leur disponibilité, connecteur et tarif. & Haute & 8\\\hline
  US-13 & En tant qu'opérateur, je veux consulter la fiche détaillée d'une borne avec
  ses connecteurs, sessions, historique et alertes. & Haute & 8\\\hline
  US-14 & En tant que plateforme, je veux recevoir et valider automatiquement les
  messages OCPP 1.6 JSON des bornes physiques ou simulées. & Haute & 13\\\hline
  US-15 & En tant que plateforme, je veux calculer l'état des bornes à partir des
  événements OCPP, transactions, heartbeats et délais de non-réponse. & Haute & 13\\\hline
  US-16 & En tant qu'opérateur, je veux superviser en temps réel les états disponible,
  occupée, non disponible, maintenance, déconnectée ou en défaut. & Haute & 13\\\hline
  US-17 & En tant qu'opérateur, je veux consulter un tableau de bord avec disponibilité,
  énergie, sessions actives et taux d'indisponibilité. & Haute & 8\\\hline
  US-18 & En tant qu'administrateur, je veux consulter le tableau de bord de mon
  organisation avec utilisateurs, clients, revenus, régions et bornes. & Haute & 8\\\hline
  US-19 & En tant qu'opérateur, je veux recevoir des alertes lorsqu'une borne est
  déconnectée, en panne, en surchauffe ou en erreur de communication. & Haute & 8\\\hline
  US-20 & En tant qu'opérateur, je veux filtrer, suivre et assigner les alertes selon
  leur gravité, statut et borne. & Moyenne & 5\\\hline
  US-21 & En tant que technicien, je veux consulter les alertes et interventions qui
  me sont assignées. & Haute & 5\\\hline
  US-22 & En tant que technicien, je veux créer un rapport d'intervention avec
  diagnostic, actions, pièces et photos. & Moyenne & 8\\\hline
  US-23 & En tant qu'opérateur, je veux planifier et suivre les maintenances
  préventives et correctives. & Moyenne & 8\\\hline
  US-24 & En tant que client, je veux démarrer une session sur une borne disponible,
  quelle que soit son organisation. & Haute & 8\\\hline
  US-25 & En tant que client, je veux arrêter ma session et consulter le résumé de
  durée, énergie et coût. & Haute & 5\\\hline
  US-26 & En tant qu'utilisateur autorisé, je veux consulter l'historique des sessions
  de mon périmètre avec durée, énergie, coût et borne. & Moyenne & 5\\\hline
  US-27 & En tant que nouvel utilisateur, je veux un onboarding adapté à mon rôle afin
  de configurer mon espace et atteindre une première action utile. & Moyenne & 5\\\hline
  US-28 & En tant que client, je veux utiliser une carte RFID ou un QR code pour
  m'identifier et lancer une recharge. & Moyenne & 8\\\hline
  US-29 & En tant que client, je veux payer mes sessions, suivre le paiement et
  prévisualiser ou télécharger mes reçus. & Moyenne & 8\\\hline
  US-30 & En tant que plateforme, je veux traiter les paiements au moyen d'un adaptateur
  simulé et remplaçable. & Moyenne & 8\\\hline
  US-31 & En tant qu'administrateur, je veux configurer les tarifs selon le kWh, la
  durée ou les plages horaires et les affecter aux actifs. & Moyenne & 8\\\hline
  US-32 & En tant que client, je veux choisir les plans des organisations et gérer
  simultanément plusieurs abonnements. & Moyenne & 8\\\hline
  US-33 & En tant qu'utilisateur autorisé, je veux consulter des analyses adaptées à
  mon rôle et échanger des rapports opérationnels. & Moyenne & 8\\\hline
  US-34 & En tant qu'utilisateur autorisé, je veux exporter les données et rapports de
  mon périmètre en PDF, JSON ou CSV. & Faible & 5\\\hline
  US-35 & En tant qu'administrateur, je veux gérer les paramètres locaux de mon
  organisation : TVA, logo, langue, fuseau et notifications. & Faible & 5\\\hline
  US-36 & En tant que super administrateur, je veux configurer les intégrations
  globales : paiement, notifications, cartographie, OAuth2, OCPP et sécurité. & Faible & 8\\\hline
  US-37 & En tant qu'administrateur, je veux gérer les contrats, garanties, notices et
  plans techniques liés aux bornes. & Faible & 5\\\hline
  US-38 & En tant qu'administrateur, je veux gérer les mises à jour firmware des bornes
  à distance. & Faible & 13\\\hline
  US-39 & En tant que visiteur, je veux consulter la présentation, contacter la
  plateforme et envoyer une demande de démonstration. & Moyenne & 3\\\hline
  US-40 & En tant qu'utilisateur authentifié, je veux recevoir des notifications
  in-app et email et gérer mes préférences. & Moyenne & 5\\\hline
  US-41 & En tant que super administrateur, je veux gérer les offres SaaS, essais,
  abonnements et factures des organisations. & Haute & 8\\\hline
  US-42 & En tant qu'administrateur, je veux suivre l'abonnement commercial, les
  limites et factures de mon organisation et demander un changement d'offre. & Moyenne & 5\\
  \hline
  \multicolumn{2}{r}{\textbf{Total}} &
  \multicolumn{2}{c}{\textbf{301 points}}\\
  \hline
\end{longtable}
\endgroup

\begin{table}[H]
  \centering
  \caption{Répartition du backlog par priorité}
  \label{tab:backlog-priority-summary}
  \begin{tabular}{@{}lrrr@{}}
    \toprule
    \textbf{Priorité} & \textbf{User stories} & \textbf{Points} &
    \textbf{Part des points}\\
    \midrule
    Haute & 20 & 160 & 53,2\,\%\\
    Moyenne & 17 & 105 & 34,9\,\%\\
    Faible & 5 & 36 & 12,0\,\%\\
    \midrule
    \textbf{Total} & \textbf{42} & \textbf{301} & \textbf{100\,\%}\\
    \bottomrule
  \end{tabular}
\end{table}

\subsection{Planification des Sprints}

Le Product Backlog est réalisé au moyen de Sprints courts, chacun associé à un
objectif vérifiable. La cadence prévue est hebdomadaire ; le dernier Sprint est
légèrement raccourci afin de respecter la date de fin du stage. Les périodes du
tableau~\ref{tab:sprint-backlogs} représentent la planification de référence. Le
contenu détaillé d'un Sprint peut être ajusté lors du Sprint Planning ou réordonné
après une Sprint Review sans modifier l'objectif global du produit.

Les story points ne sont pas convertis en heures. Ils servent à comparer l'effort et
le risque des user stories. La capacité est appréciée progressivement à partir des
Sprints précédents, ce qui évite de présenter une précision artificielle au début du
stage.

\begingroup
\small
\setlength{\tabcolsep}{2.5pt}
\begin{longtable}{|C{1.4cm}|C{2.25cm}|C{2.6cm}|p{3.3cm}|p{4.4cm}|}
  \caption{Synthèse des Sprint Backlogs}
  \label{tab:sprint-backlogs}\\
  \hline
  \textbf{Sprint} & \textbf{Période prévue} & \textbf{Objectif} &
  \textbf{User stories principales} & \textbf{Résultat inspecté en revue}\\
  \hline
  \endfirsthead
  \hline
  \textbf{Sprint} & \textbf{Période prévue} & \textbf{Objectif} &
  \textbf{User stories principales} & \textbf{Résultat inspecté en revue}\\
  \hline
  \endhead
  S01 & 1--7 juillet &
  Cadrer le produit et valider l'expérience cible &
  US-27, US-39 et travaux de conception &
  Backlog initial, acteurs, cas d'utilisation global, architecture de référence et
  prototype navigable.\\\hline
  S02 & 8--14 juillet &
  Sécuriser l'entrée et l'organisation des comptes &
  US-01 à US-06 &
  Authentification, création de compte, invitation, activation, rôles et isolation
  d'une organisation.\\\hline
  S03 & 15--21 juillet &
  Constituer et localiser le réseau de recharge &
  US-07 à US-13, US-37 &
  Stations, connecteurs, cartographie, fiche détaillée et documents accessibles selon
  le rôle.\\\hline
  S04 & 22--28 juillet &
  Produire une disponibilité issue du terrain &
  US-14 à US-16, US-19 &
  Connexion d'une station simulée, réception des événements, projection des états,
  historique et alerte de connectivité.\\\hline
  S05 & 29 juillet--4 août &
  Réaliser une recharge complète &
  US-24 à US-30 &
  Sélection d'une station, connexion, autorisation, mesures, arrêt, paiement simulé
  et reçu.\\\hline
  S06 & 5--11 août &
  Traiter un incident jusqu'à sa résolution &
  US-20 à US-23, US-40 &
  Alerte qualifiée, intervention assignée, maintenance, preuves de travail et
  notifications.\\\hline
  S07 & 12--18 août &
  Structurer les offres et la tarification &
  US-31, US-32, US-41, US-42 &
  Tarifs de recharge, adhésions des clients, essais et abonnements commerciaux des
  organisations.\\\hline
  S08 & 19--25 août &
  Fournir les outils de pilotage &
  US-05, US-17, US-18, US-33 à US-36 &
  Tableaux de bord spécialisés, rapports, exports, audit, paramètres et état des
  intégrations.\\\hline
  S09 & 26--31 août &
  Stabiliser et préparer la recette &
  Tests transversaux et documentation &
  Corrections finales, sécurité, tests de recette, API documentée, manuels, rapport
  et préparation du déploiement.\\
  \hline
\end{longtable}
\endgroup

La gestion distante du firmware (US-38) reste volontairement dans le Product Backlog
sans être engagée dans un Sprint du MVP. Elle nécessite une validation matérielle et
sécuritaire qui dépasse le parcours principal du stage. Cette transparence évite de
présenter comme réalisée une évolution seulement envisagée.

\subsection{Matrice de traçabilité fonctionnelle}

La matrice suivante relie les grands groupes de besoins à leurs principaux contrats
API et écrans. Les routes détaillées et la spécification OpenAPI seront présentées
dans les livrables techniques.

\begingroup
\small
\setlength{\tabcolsep}{3pt}
\begin{longtable}{|p{2.9cm}|p{2.9cm}|p{4.2cm}|p{3.9cm}|}
  \caption{Traçabilité entre backlog, API et interfaces}
  \label{tab:functional-traceability}\\
  \hline
  \textbf{Domaine} & \textbf{User stories} & \textbf{Contrats API principaux} &
  \textbf{Interfaces principales}\\
  \hline
  \endfirsthead
  \hline
  \textbf{Domaine} & \textbf{User stories} & \textbf{Contrats API principaux} &
  \textbf{Interfaces principales}\\
  \hline
  \endhead
  Identité et entrée & US-01--06, 27, 39 &
  Auth, profil, onboarding, utilisateurs, organisations, demandes de démo &
  Landing, connexion, activation, bienvenue, profil, utilisateurs, organisations\\\hline
  Stations et carte & US-07--13, 37 &
  Stations, commissioning, connecteurs, carte, télémétrie, documents &
  Stations, recherche de station, détail station, cartes par rôle\\\hline
  OCPP et disponibilité & US-14--16, 19 &
  API interne OCPP, commandes, transitions, alertes &
  Détail station, supervision, carte, centre d'alertes\\\hline
  Exploitation & US-20--23, 40 &
  Alertes, interventions, maintenances, notifications &
  Alertes, interventions, maintenance, calendrier, cloche de notifications\\\hline
  Recharge et paiement & US-24--30 &
  Tentatives, sessions, paiements, reçu, webhook fournisseur &
  Assistant de recharge, sessions, vue live, paiements et reçus\\\hline
  Tarification et commerce & US-31, 32, 41, 42 &
  Tarifs, plans, abonnements clients, portefeuille commercial, factures &
  Tarifs, abonnements, portefeuille commercial, facturation organisationnelle\\\hline
  Pilotage et rapports & US-05, 17--18,\newline 33--36 &
  Dashboards, reporting, rapports internes, exports, audit, intégrations, paramètres &
  Tableaux de bord par rôle, analyses, rapports, audit, intégrations, paramètres\\
  \hline
\end{longtable}
\endgroup
